\lysh{16147_SOLUTION}{1}
ΛΥΣΗ
% TODO: TikZ σχήμα
Τα διανύσματα που δίνονται έχουν συντεταγμένες $\vec{\alpha}=(3, 3\sqrt{3})$, $\vec{\beta}=(\sqrt{2}, 0)$, $\vec{\gamma}=(0, -3)$ και $\vec{\delta}=(-1, 1)$.

\noindent
α) Ο συντελεστής διεύθυνσης διανύσματος, όταν η τετμημένη του δεν είναι μηδέν, ορίζεται ως το πηλίκο τεταγμένης του διανύσματος προς τετμημένη του διανύσματος. Οπότε
$\lambda_{\vec{\alpha}} = \frac{3\sqrt{3}}{3} = \sqrt{3}$, $\lambda_{\vec{\beta}} = \frac{0}{\sqrt{2}} = 0$, $\lambda_{\vec{\gamma}}$ δεν ορίζεται, $\lambda_{\vec{\delta}} = \frac{1}{-1} = -1$.

\noindent
β) Γνωρίζουμε ότι η εφαπτομένη της γωνίας που σχηματίζει ένα διάνυσμα με το θετικό ημιάξονα $Ox$ ισούται με το συντελεστή διεύθυνσης του διανύσματος.
Αν $\omega$ είναι η γωνία που σχηματίζει το διάνυσμα $\vec{\alpha}$ με το θετικό ημιάξονα $Ox$, επειδή $\lambda_{\vec{\alpha}} = \sqrt{3}$ (από ερώτημα α), θα ισχύει $\varepsilon\varphi\omega = \sqrt{3}$. Επιπλέον το πέρας του διανύσματος βρίσκεται στο $1^ο$ τεταρτημόριο, αφού έχει θετικές συντεταγμένες, άρα $\omega = 60^\circ$.
Το διάνυσμα $\vec{\beta} = (\sqrt{2}, 0)$, έχει τεταγμένη $0$, άρα $\vec{\beta} // x'x$, και επειδή το $\vec{\beta}$ έχει θετική τετμημένη σχηματίζει με το θετικό ημιάξονα $Ox$ γωνία $0^\circ$.
Το διάνυσμα $\vec{\gamma} = (0, -3)$, έχει τετμημένη $0$, άρα $\vec{\gamma} // y'y$, και επειδή το $\vec{\gamma}$ έχει αρνητική τεταγμένη σχηματίζει με το θετικό ημιάξονα $Ox$ γωνία $270^\circ$.
Αν $\varphi$ είναι η γωνία που σχηματίζει το διάνυσμα $\vec{\delta}$ με το θετικό ημιάξονα $Ox$, επειδή $\lambda_{\vec{\delta}} = -1$ (από ερώτημα α), θα ισχύει $\varepsilon\varphi\varphi = -1$. Επιπλέον το διάνυσμα έχει αρνητική τετμημένη και θετική τεταγμένη, οπότε το πέρας του βρίσκεται στο $2^ο$ τεταρτημόριο, άρα $\varphi = 145^\circ$.

\noindent
γ)
\[
|\vec{\alpha}| = \sqrt{3^2 + (3\sqrt{3})^2} = \sqrt{9 + 27} = \sqrt{36} = 6
\]
\[
|\vec{\gamma}| = \sqrt{0^2 + (-3)^2} = \sqrt{0 + 9} = \sqrt{9} = 3.
\]
\endlysh